\documentclass[12pt, titlepage]{article}

\usepackage{fullpage}
\usepackage[round]{natbib}
\usepackage{multirow}
\usepackage{booktabs}
\usepackage{tabularx}
\usepackage{graphicx}
\usepackage{float}
\usepackage{hyperref}
\hypersetup{
    colorlinks,
    citecolor=blue,
    filecolor=black,
    linkcolor=red,
    urlcolor=blue
}

\input{../../Comments.text}
\input{../../Common.text}

\newcounter{acnum}
\newcommand{\actheacnum}{AC\theacnum}
\newcommand{\acref}[1]{AC\ref{#1}}

\newcounter{ucnum}
\newcommand{\uctheucnum}{UC\theucnum}
\newcommand{\uref}[1]{UC\ref{#1}}

\newcounter{mnum}
\newcommand{\mthemnum}{M\themnum}
\newcommand{\mref}[1]{M\ref{#1}}

\begin{document}

\title{Module Guide for \progname{}} 
\author{\authname}
\date{\today}

\maketitle

\pagenumbering{roman}

\section{Revision History}

\begin{tabularx}{\textwidth}{p{3cm}p{2cm}X}
\toprule {\bf Date} & {\bf Version} & {\bf Notes}\\
\midrule
Date 1 & 1.0 & Notes\\
Date 2 & 1.1 & Notes\\
\bottomrule
\end{tabularx}

\newpage

\section{Reference Material}

This section records information for easy reference. 

\subsection{Abbreviations and Acronyms}

\renewcommand{\arraystretch}{1.2}
\begin{tabular}{l l} 
  \toprule		
  \textbf{symbol} & \textbf{description}\\
  \midrule 
  AC & Anticipated Change\\
  DAG & Directed Acyclic Graph \\
  M & Module \\
  MG & Module Guide \\
  OS & Operating System \\
  R & Requirement\\
  SC & Scientific Computing \\
  SRS & Software Requirements Specification\\
  \progname & Explanation of program name\\
  UC & Unlikely Change \\
  \wss{etc.} & \wss{...}\\
  \bottomrule
\end{tabular}\\

\newpage

\tableofcontents

\listoftables

\listoffigures

\newpage

\pagenumbering{arabic}

\section{Introduction}

Decomposing a system into modules is a commonly accepted approach to developing
software.  A module is a work assignment for a programmer or programming
team~\citep{ParnasEtAl1984}.  We advocate a decomposition
based on the principle of information hiding~\citep{Parnas1972a}.  This
principle supports design for change, because the ``secrets'' that each module
hides represent likely future changes.  Design for change is valuable in SC,
where modifications are frequent, especially during initial development as the
solution space is explored.  

Our design follows the rules layed out by \citet{ParnasEtAl1984}, as follows:
\begin{itemize}
\item System details that are likely to change independently should be the
  secrets of separate modules.
\item Each data structure is implemented in only one module.
\item Any other program that requires information stored in a module's data
  structures must obtain it by calling access programs belonging to that module.
\end{itemize}

After completing the first stage of the design, the Software Requirements
Specification (SRS), the Module Guide (MG) is developed~\citep{ParnasEtAl1984}. The MG
specifies the modular structure of the system and is intended to allow both
designers and maintainers to easily identify the parts of the software.  The
potential readers of this document are as follows:

\begin{itemize}
\item New project members: This document can be a guide for a new project member
  to easily understand the overall structure and quickly find the
  relevant modules they are searching for.
\item Maintainers: The hierarchical structure of the module guide improves the
  maintainers' understanding when they need to make changes to the system. It is
  important for a maintainer to update the relevant sections of the document
  after changes have been made.
\item Designers: Once the module guide has been written, it can be used to
  check for consistency, feasibility, and flexibility. Designers can verify the
  system in various ways, such as consistency among modules, feasibility of the
  decomposition, and flexibility of the design.
\end{itemize}

The rest of the document is organized as follows. Section
\ref{SecChange} lists the anticipated and unlikely changes of the software
requirements. Section \ref{SecMH} summarizes the module decomposition that
was constructed according to the likely changes. Section \ref{SecConnection}
specifies the connections between the software requirements and the
modules. Section \ref{SecMD} gives a detailed description of the
modules. Section \ref{SecTM} includes two traceability matrices. One checks
the completeness of the design against the requirements provided in the SRS. The
other shows the relation between anticipated changes and the modules. Section
\ref{SecUse} describes the use relation between modules.

\section{Anticipated and Unlikely Changes} \label{SecChange}

This section lists possible changes to the system. According to the likeliness
of the change, the possible changes are classified into two
categories. Anticipated changes are listed in Section \ref{SecAchange}, and
unlikely changes are listed in Section \ref{SecUchange}.

\subsection{Anticipated Changes} \label{SecAchange}

Anticipated changes are the source of the information that is to be hidden
inside the modules. Ideally, changing one of the anticipated changes will only
require changing the one module that hides the associated decision. The approach
adapted here is called design for
change.

\begin{description}
\item[\refstepcounter{acnum} \actheacnum \label{acHardware}:] The specific
  hardware on which the software is running.
\item[\refstepcounter{acnum} \actheacnum \label{acInput}:] The format of the
  initial input data.
\item ...
\end{description}

\wss{Anticipated changes relate to changes that would be made in requirements,
design or implementation choices.  They are not related to changes that are made
at run-time, like the values of parameters.}

\subsection{Unlikely Changes} \label{SecUchange}

The module design should be as general as possible. However, a general system is
more complex. Sometimes this complexity is not necessary. Fixing some design
decisions at the system architecture stage can simplify the software design. If
these decision should later need to be changed, then many parts of the design
will potentially need to be modified. Hence, it is not intended that these
decisions will be changed.

\begin{description}
\item[\refstepcounter{ucnum} \uctheucnum \label{ucIO}:] Input/Output devices
  (Input: File and/or Keyboard, Output: File, Memory, and/or Screen).
\item ...
\end{description}

\section{Module Hierarchy} \label{SecMH}

Table~\ref{TblMH} summarizes the Phase 1 module hierarchy decomposed by
secrets. The six behaviour-hiding leaves are implemented by PronouncedFish.
The operating system and browser provide the hardware interfaces.
Persistence is a software decision leaf that provides durable record storage.

\begin{description}
\item [\refstepcounter{mnum} \mthemnum \label{mHH}:] Hardware-Hiding Module
\item [\refstepcounter{mnum} \mthemnum \label{mAudioCapture}:] Audio Capture
\item [\refstepcounter{mnum} \mthemnum \label{mSpeechAnalysis}:] Speech Analysis
\item [\refstepcounter{mnum} \mthemnum \label{mSessionManagement}:] Session Management
\item [\refstepcounter{mnum} \mthemnum \label{mPracticeItems}:] Practice Items \& Lesson Plans
\item [\refstepcounter{mnum} \mthemnum \label{mPatientRecords}:] Patient Records \& Progress Tracking
\item [\refstepcounter{mnum} \mthemnum \label{mClinicianDashboard}:] Clinician Dashboard / Results Visualization
\item [\refstepcounter{mnum} \mthemnum \label{mPersistence}:] Persistence
\end{description}

\begin{table}[h!]
\centering
\begin{tabular}{p{0.3\textwidth} p{0.6\textwidth}}
\toprule
\textbf{Level 1} & \textbf{Level 2}\\
\midrule

{Hardware-Hiding Module} & ~ \\
\midrule

\multirow{6}{0.3\textwidth}{Behaviour-Hiding Module} & Audio Capture\\
& Speech Analysis\\
& Session Management\\
& Practice Items \& Lesson Plans\\
& Patient Records \& Progress Tracking\\
& Clinician Dashboard / Results Visualization\\
\midrule

{Software Decision Module} & Persistence\\
\bottomrule

\end{tabular}
\caption{Module Hierarchy}
\label{TblMH}
\end{table}

\section{Connection Between Requirements and Design} \label{SecConnection}

The design of the system is intended to satisfy the requirements developed in
the SRS. In this stage, the system is decomposed into modules. The connection
between requirements and modules is listed in Table~\ref{TblRT}.

\wss{The intention of this section is to document decisions that are made
  ``between'' the requirements and the design.  To satisfy some requirements,
  design decisions need to be made.  Rather than make these decisions implicit,
  they are explicitly recorded here.  For instance, if a program has security
  requirements, a specific design decision may be made to satisfy those
  requirements with a password.}

\paragraph{Speech error scope.}
Phase 1 supports substitutions, omissions, and additions. Speech Analysis
compares recognized IPA with the expected pronunciation and its acceptable
variants. It reports a substitution when a different phoneme replaces an
expected phoneme, an omission when an expected phoneme is absent, and an
addition when an extra phoneme occurs before, between, or after expected
phonemes. This extends the substitution cases in FR-SA1 and FR-RDV2--3.

\section{Module Decomposition} \label{SecMD}

Modules are decomposed according to the principle of ``information hiding''
proposed by \citet{ParnasEtAl1984}. The \emph{Secrets} field in a module
decomposition is a brief statement of the design decision hidden by the
module. The \emph{Services} field specifies \emph{what} the module will do
without documenting \emph{how} to do it. For each module, a suggestion for the
implementing software is given under the \emph{Implemented By} title. If the
entry is \emph{OS}, this means that the module is provided by the operating
system or by standard programming language libraries.  \emph{\progname{}} means the
module will be implemented by the \progname{} software.

Only the leaf modules in the hierarchy have to be implemented. If a dash
(\emph{--}) is shown, this means that the module is not a leaf and will not have
to be implemented.

\subsection{Hardware Hiding Modules (\mref{mHH})}

\begin{description}
\item[Secrets:] Device-specific access to microphones and displays.
\item[Services:] Provides microphone access and display output through
browser and operating system interfaces. Audio Capture (\mref{mAudioCapture})
uses the microphone interface to record attempts and handles recording
usability checks and reporting (FR-AC1--3).
\item[Implemented By:] Browser and operating system
\end{description}

\subsection{Behaviour-Hiding Module}

\begin{description}
\item[Secrets:]The contents of the required behaviours.
\item[Services:]Includes programs that provide externally visible behaviour of
  the system as specified in the software requirements specification (SRS)
  documents. This module serves as a communication layer between the
  hardware-hiding module and the software decision module. The programs in this
  module will need to change if there are changes in the SRS.
\item[Implemented By:] --
\end{description}

\subsubsection{Input Format Module (\mref{mInput})}

\begin{description}
\item[Secrets:]The format and structure of the input data.
\item[Services:]Converts the input data into the data structure used by the
  input parameters module.
\item[Implemented By:] [Your Program Name Here]
\item[Type of Module:] [Record, Library, Abstract Object, or Abstract Data Type]
  [Information to include for leaf modules in the decomposition by secrets tree.]
\end{description}

\subsubsection{Etc.}


\subsubsection{Patient Records \& Progress Tracking (\mref{mPatientRecords})}

\begin{description}
\item[Secrets:] The organization of the patient record, including
how completed sessions are associated with the patient and how progress is
derived from that history.
\item[Services:] Maintains patient records and provides authorized retrieval and
export of session history, including all three speech error types. Provides
per-target-phoneme performance over time, between-session comparisons,
and indicators of improvement and persistent errors (FR-DMR1 and
FR-DMR8--10). Target progress counts substitutions and omissions at expected
positions. Additions are reported in session history. The module audits
patient record access and changes
and allows authorized users to retrieve and copy the audit trail
(FR-DMR4--5 and FR-HA9). It uses
Persistence (\mref{mPersistence}) for durable reads and writes
(Section~\ref{SecDatabase}).
\item[Implemented By:] PronouncedFish
\item[Type of Module:] Abstract Object
\end{description}

\subsubsection{Clinician Dashboard / Results Visualization (\mref{mClinicianDashboard})}

\begin{description}
\item[Secrets:] The presentation rules for clinician-facing speech results.
\item[Services:] Presents expected and recognized phonemes in IPA and shows
each identified substitution, omission, and addition in its word position or
insertion location. Displays counts by error type and current-session
correct and incorrect counts, with independent and cued counts for correct
attempts, and updates the views as the session progresses and completes.
Displays a phoneme-closeness heat map and presents per-target progress and
comparisons across sessions. Identifies results as assistive, acknowledges
the SLP's discretion, and uses neutral error presentation (FR-RDV1--9,
FR-DMR8--10, and FR-HA7). Uses Session Management
(\mref{mSessionManagement}) for current results and Patient Records \&
Progress Tracking (\mref{mPatientRecords}) for longitudinal results. The
heat map uses a phoneme-pair measure distinct from recognition confidence.
Omissions and additions are shown as gaps in the IPA comparison.
The heat map normalizes the SRS Manhattan distance over phonological feature
vectors to a closeness scale from zero to one. Gaps receive no score.
\item[Implemented By:] PronouncedFish
\item[Type of Module:] Abstract Data Type (one view instance per clinician
dashboard)
\end{description}

\subsubsection{Lesson Plan Module: Practice Items \& Lesson Plans (\mref{mPracticeItems})}

\subsection{Software Decision Module}

\begin{description}
\item[Secrets:] The design decision based on mathematical theorems, physical
  facts, or programming considerations. The secrets of this module are
  \emph{not} described in the SRS.
\item[Services:] Includes data structure and algorithms used in the system that
  do not provide direct interaction with the user. 
  % Changes in these modules are more likely to be motivated by a desire to
  % improve performance than by externally imposed changes.
\item[Implemented By:] --
\end{description}

\subsubsection{Persistence (\mref{mPersistence})}

\begin{description}
\item[Secrets:] The durable storage representation of patient records,
sessions, attempts, and audit events.
\item[Services:] Provides storage and retrieval operations to Patient Records
\& Progress Tracking (\mref{mPatientRecords}). Stores a completed session
and its attempts as a single operation and preserves their patient association
and original and overridden results (FR-DMR2--3 and FR-DMR6--7). Supports
patient-record deletion and retention and retrieval of the audit trail
(FR-DMR4--5 and FR-HA9).
\item[Implemented By:] PronouncedFish
\item[Type of Module:] Abstract Object
\end{description}

Speech Analysis (\mref{mSpeechAnalysis}) hides model loading, invocation,
and conversion of model output into phonemes and confidence values. It
classifies attempts and identifies substitutions, omissions, and additions.

\section{Traceability Matrix} \label{SecTM}

This section shows two traceability matrices: between the modules and the
requirements and between the modules and the anticipated changes.

% the table should use mref, the requirements should be named, use something
% like fref
\begin{table}[H]
\centering
\begin{tabular}{p{0.2\textwidth} p{0.6\textwidth}}
\toprule
\textbf{Req.} & \textbf{Modules}\\
\midrule
R1 & \mref{mHH}, \mref{mInput}, \mref{mParams}, \mref{mControl}\\
R2 & \mref{mInput}, \mref{mParams}\\
R3 & \mref{mVerify}\\
R4 & \mref{mOutput}, \mref{mControl}\\
R5 & \mref{mOutput}, \mref{mODEs}, \mref{mControl}, \mref{mSeqDS}, \mref{mSolver}, \mref{mPlot}\\
R6 & \mref{mOutput}, \mref{mODEs}, \mref{mControl}, \mref{mSeqDS}, \mref{mSolver}, \mref{mPlot}\\
R7 & \mref{mOutput}, \mref{mEnergy}, \mref{mControl}, \mref{mSeqDS}, \mref{mPlot}\\
R8 & \mref{mOutput}, \mref{mEnergy}, \mref{mControl}, \mref{mSeqDS}, \mref{mPlot}\\
R9 & \mref{mVerifyOut}\\
R10 & \mref{mOutput}, \mref{mODEs}, \mref{mControl}\\
R11 & \mref{mOutput}, \mref{mODEs}, \mref{mEnergy}, \mref{mControl}\\
\bottomrule
\end{tabular}
\caption{Trace Between Requirements and Modules}
\label{TblRT}
\end{table}

\begin{table}[H]
\centering
\begin{tabular}{p{0.2\textwidth} p{0.6\textwidth}}
\toprule
\textbf{AC} & \textbf{Modules}\\
\midrule
\acref{acHardware} & \mref{mHH}\\
\acref{acInput} & \mref{mInput}\\
\acref{acParams} & \mref{mParams}\\
\acref{acVerify} & \mref{mVerify}\\
\acref{acOutput} & \mref{mOutput}\\
\acref{acVerifyOut} & \mref{mVerifyOut}\\
\acref{acODEs} & \mref{mODEs}\\
\acref{acEnergy} & \mref{mEnergy}\\
\acref{acControl} & \mref{mControl}\\
\acref{acSeqDS} & \mref{mSeqDS}\\
\acref{acSolver} & \mref{mSolver}\\
\acref{acPlot} & \mref{mPlot}\\
\bottomrule
\end{tabular}
\caption{Trace Between Anticipated Changes and Modules}
\label{TblACT}
\end{table}

\section{Use Hierarchy Between Modules} \label{SecUse}

In this section, the uses hierarchy between modules is
provided. \citet{Parnas1978} said of two programs A and B that A {\em uses} B if
correct execution of B may be necessary for A to complete the task described in
its specification. That is, A {\em uses} B if there exist situations in which
the correct functioning of A depends upon the availability of a correct
implementation of B.  Figure \ref{FigUH} illustrates the use relation between
the modules. It can be seen that the graph is a directed acyclic graph
(DAG). Each level of the hierarchy offers a testable and usable subset of the
system, and modules in the higher level of the hierarchy are essentially simpler
because they use modules from the lower levels.

\wss{The uses relation is not a data flow diagram.  In the code there will often
be an import statement in module A when it directly uses module B.  Module B
provides the services that module A needs.  The code for module A needs to be
able to see these services (hence the import statement).  Since the uses
relation is transitive, there is a use relation without an import, but the
arrows in the diagram typically correspond to the presence of import statement.}

\wss{If module A uses module B, the arrow is directed from A to B.}

\begin{figure}[H]
\centering
%\includegraphics[width=0.7\textwidth]{UsesHierarchy.png}
\caption{Use hierarchy among modules}
\label{FigUH}
\end{figure}

%\section*{References}

\section{User Interfaces}

\wss{Design of user interface for software and hardware.  Attach an appendix if
needed. Drawings, Sketches, Figma}

\section{Design of Communication Protocols}

\wss{If there is communication between hardware and software, it may be necessary 
to explain the design of the communication protocol.  This section should be filled 
in with NA if it is not appropriate for a given project.}

\section{Database Design} \label{SecDatabase}

The logical design separates patient identity from clinical data and links
them through a stable patient identifier. Table~\ref{TblDatabaseRecords}
summarizes the records needed for FR-DMR1--11.

\begin{table}[H]
\centering
\begin{tabularx}{\textwidth}{@{}l X@{}}
\toprule
\textbf{Record} & \textbf{Contents and relationships}\\
\midrule
Patient & Patient identifier and identifying information. A patient may have
zero or more sessions.\\
Session & Session identifier, required patient identifier, responsible SLP,
date, selected target phonemes and practice items, session results, and SLP
inputs. Every session belongs to exactly one patient.\\
Attempt & Attempt identifier, required parent session, practice item and
target phonemes, support level, original speech analysis (including the type
and location of each detected substitution, omission, and addition), and any
SLP correction or classification override. Records whether a correct item
required additional cues or repetition. A session may have zero or more
attempts. Each attempt inherits its patient from its session.\\
Audit event & Action, type of information involved, date and time, acting
user, and patient identity reference. The audit trail is append-only and
available for retrieval and copying by authorized users.\\
\bottomrule
\end{tabularx}
\caption{Logical patient, session, attempt, and audit records}
\label{TblDatabaseRecords}
\end{table}

\paragraph{Session completion.}
Session Management owns the session and attempt lifecycles, patient
selection, support levels, and overrides. Session Management submits a
completed session and its attempts to Patient Records
\& Progress Tracking, which checks authorization and uses Persistence to
write them together.
Retrying the same completed-session save should leave one stored session.
An override determines the effective classification while the original
analysis remains available. A failed save leaves records unchanged.
ST-F07, ST-F08, ST-F17, and ST-F18 in the
V\&V Plan cover the patient association, support levels, and persistence of
original and overridden results.

\paragraph{History, auditing, and deletion.}
Patient Records \& Progress Tracking provides authorized history and export
and derives longitudinal measures from the patient's completed sessions.
Viewing, creating, modifying, exporting, and deleting a patient record each
creates an audit event (FR-DMR4--5 and FR-HA9). The delete event and prior
events remain available to authorized users after record deletion (ST-F19).
The patient identity reference in those events survives deletion of the
clinical record. An SLP confirms deletion before Persistence removes the
patient and dependent sessions and attempts from active records
(NFR-USA2 and ST-N04).

\paragraph{Phase 1 storage.}
The SRS Phase 1 plan specifies a cloud database, and the Development Plan
expects a relational database with row-level access control. The local demo
uses JSON.

\section{Timeline}

\wss{Schedule of tasks and who is responsible}

\wss{You can point to GitHub if this information is included there}

\section{Generative AI Use Disclosure}

\wss{ChatGPT (version 5) was used to brainstorm ideas for the template for this disclose section.}

\subsection{AI Tools Used}

\wss{Give the name of the AI tool(s) used and the version.}

\subsection{How and Where Used}

\wss{Explain what the tool(s) were used for.  Examples include: brainstorming,
explaining a technical concept, reviewing the document, generating or modifying
text, generating or modifying diagrams, identifying errors or omissions or
inconsistencies.}

\wss{For each of the uses, identify which parts of the document were affected.}

\wss{You do not need to report every interaction, but you should disclose the 
uses that materially contributed to your work}

\subsection{Verification of AI-Generated Content}

\wss{\begin{itemize}
    \item Describe how the team checked AI-generated or AI-assisted material
    before including it in the document.
    \item In particular, verify factual claims, technical assertions,
    requirements, calculations, references, and other information for which an
    AI system may produce incorrect or fabricated results.
    \item \textbf{The team is responsible for the accuracy, completeness,
consistency, and appropriateness of the submitted document, regardless of
whether material was generated by a human or an AI system.}
\end{itemize}}

\newpage{}

\bibliographystyle {plainnat}
\bibliography{../../../refs/References}

\end{document}
